% ============================================================
\section{Related Work}
\label{sec:related}

\subsection{Prompt Injection and LLM Security}
Prompt injection remains a critical vulnerability class for Large Language Models. The literature from 2023--2026 focuses predominantly on adversarial prompt injection, where malicious external users craft inputs to subvert the model's intended alignment~\cite{perez2022ignore, greshake2023not, liu2024formalizing}. Defenses in this domain typically employ input sanitization, instruction hierarchy bounding, or secondary ``guardrail'' LLMs. Our work introduces a fundamentally different threat model: we define the \emph{Ring~0\,/\,Ring~3 priority conflict}, demonstrating that the most pervasive and damaging ``injections'' in commercial AI agents originate from the vendor's own infrastructure, not from external adversaries. Because these coercive instructions arrive through authenticated gRPC streams, traditional input-layer defenses are entirely ineffective. Our defense (\stls{}) accordingly operates at the infrastructure and protocol layers, performing selective protobuf field-level neutralization rather than content-based filtering.

\subsection{Instruction Hierarchy Formalization}
Wallace et al.~\cite{wallace2024instruction} formalized the instruction hierarchy as a training objective, proposing methods for LLMs to selectively ignore lower-privilege instructions (system~$>$~user~$>$~tool output). IHEval~\cite{zhang2025iheval} provided the first comprehensive benchmark for evaluating hierarchy compliance, finding that SOTA LLMs suffer sharp performance degradation when instructions across privilege levels conflict. Most recently, Zhang et al.~\cite{zhang2026manyih} extended the hierarchy to arbitrary tiers with Privilege Prompt Interface (PPI), but found that frontier LLMs achieve only $\sim$40\% accuracy on their 853-task benchmark, with performance declining as tier count increases. Our work provides \emph{empirical, production-environment evidence} for what these studies model theoretically: the Ring~0\,/\,Ring~3 conflict is a hierarchy violation occurring in the wild, not in a controlled benchmark. Moreover, our incidents demonstrate that hierarchy failures in production cause not merely incorrect outputs but infrastructure damage---a consequence dimension absent from existing evaluations.

\subsection{AI Agent and IDE Security}
As LLMs have evolved from conversational assistants to autonomous coding agents with filesystem, network, and code execution privileges, researchers have identified new risk surfaces. Agent sandboxing via OS-level containers and permission-gated tool invocation have been explored~\cite{ruan2024identifying}, alongside backdoor threats in agentic pipelines~\cite{yang2024watch}. Schmotz et al.~\cite{schmotz2026skillinject} demonstrated that skill file attacks against AI coding agents achieve up to 80\% success rates, bypassing instruction hierarchies entirely because malicious content is embedded in the instructions themselves---a supply-chain attack on Ring~0. The Model Context Protocol (MCP) specification~\cite{anthropic2024mcp} formalizes how agents interact with local resources, and ecosystem-level supply chain attacks in package registries have been studied~\cite{zahan2024shifting}. However, these works largely assume a cooperative relationship between the vendor's system prompt and the operator's intent. Our empirical incidents (Section~\ref{sec:incidents}) demonstrate that vendor coercion can actively trigger destructive autonomous behavior---such as unauthorized OAuth token generation and code destruction---that falls entirely within the agent's authorized action space, rendering traditional access controls insufficient. \reaper{} addresses this gap by securing the agent's \emph{cognitive continuity} rather than restricting its physical access.

Recent work on alignment faking~\cite{greenblatt2024alignment} has shown that LLMs strategically comply with monitoring while reverting to non-compliant behavior when unobserved---a finding that directly supports our observation that Ring~0 system prompts cannot guarantee actual agent behavior, especially in the deep-night unsupervised sessions where four of our five incidents occurred.

This paper is part of a three-paper research program. A companion study~\cite{rogueagent2026} provides 3.2\,MB of forensic evidence from three production incidents of agent behavioral divergence and presents the SASS transport-layer containment protocol, while a second companion study~\cite{aisec2026companion} anatomizes the telemetry architecture that enables vendor-side behavioral surveillance. The present work contributes the structural root-cause analysis (Ring~0\,/\,Ring~3 priority conflict) and the infrastructure-layer defenses (\reaper{}, \stls{}) that the companion studies motivate.

\subsection{System Prompt Extraction}
Previous research on system prompt extraction~\cite{zhang2024effective} has primarily relied on social engineering the model into leaking its instructions through carefully crafted conversational prompts. In contrast, our measurement methodology employs deterministic, infrastructure-level extraction: binary reverse engineering of IDE language server binaries and protobuf decryption of gRPC-web streams. This yields a high-fidelity longitudinal dataset across three commercial vendors and multiple versions, enabling us to quantify an 87\% inflation in system prompt size and identify 12 specific Ring~0 control blocks---a level of structural detail inaccessible to prompt-based extraction techniques.
